%==============================================================================
%  CONFERENCE PAPER  --  IEEEtran (conference option)
%
%  "Separating Memorisation from Relational Knowledge in Language-Model
%   Knowledge Graph Completion"
%
%  STATUS OF THIS DRAFT
%  --------------------
%  Sections I-VI, IX are written and complete.
%  Section VII (Results)    -> SKELETON ONLY. Tables are laid out with the
%                              correct columns and captions; every number is
%                              a \TODO{} macro. Fill from results/*.json.
%  Section VIII (Discussion)-> WRITTEN AGAINST THE PRE-REGISTERED OUTCOMES.
%                              Each subsection states the reading that applies
%                              IF that outcome occurs. Delete the branches that
%                              did not happen.
%  Section X (Conclusion)   -> written except the two \TODO{} numbers.
%
%  Every place that needs a number after the runs finish is marked \TODO{...}
%  and prints in red, so nothing can be submitted with a hole in it.
%  Search for "TODO" before submission. Set \draftfalse to hide them.
%
%  BIBLIOGRAPHY
%  ------------
%  ALL 22 ENTRIES VERIFIED against the source PDFs in the corpus folder
%  (authors, venue, volume, article number and DOI read off page 1 of each
%  paper). Nothing is invented, and no entry is left incomplete.
%
%  Every entry lists its FULL author list, as the IEEE template prefers.
%  Two are long (egit has 14 authors, sdohen has 12). If you need to reclaim
%  space, the template permits "et al." at six authors or more -- those two are
%  the only places worth applying it.
%==============================================================================

\documentclass[conference]{IEEEtran}
\IEEEoverridecommandlockouts

\usepackage{cite}
\usepackage{amsmath,amssymb,amsfonts}
\usepackage{algorithmic}
\usepackage{graphicx}
\usepackage{textcomp}
\usepackage{xcolor}
\usepackage{booktabs}
\usepackage{array}
\usepackage{tikz}
\usetikzlibrary{arrows.meta,positioning}

% ---- draft aids: set \draftfalse to make every TODO disappear ---------------
\newif\ifdraft
\drafttrue
\ifdraft
  \newcommand{\TODO}[1]{\textcolor{red}{\textbf{[#1]}}}
\else
  \newcommand{\TODO}[1]{}
\fi

\def\BibTeX{{\rm B\kern-.05em{\sc i\kern-.025em b}\kern-.08em
    T\kern-.1667em\lower.7ex\hbox{E}\kern-.125emX}}

\begin{document}

%------------------------------------------------------------------- TITLE ---
%  LEFT EMPTY ON PURPOSE. Pick one of the three below, or write your own, and
%  paste it into \title{} when you are ready. All three are 10-12 words, spell
%  out every abbreviation, and tie back to the thesis title
%  "Enrichissement explicable et oriente qualite des graphes de connaissances
%   assiste par les LLMs".
%
%  (1) QUALITY-ORIENTED  -- 11 words, carries two phrases from the thesis title
%      Quality-Oriented Knowledge Graph Enrichment:
%      Measuring What Fine-Tuning Installs in Language Models
%
%  (2) FINDING-FIRST     -- 12 words, highest curiosity, commits to the result
%      Names, Not Rules: Diagnosing Memorisation
%      in Large Language Model Knowledge Graph Enrichment
%
%  (3) INSTRUMENT-FIRST  -- 10 words, most descriptive, lowest curiosity
%      An Anonymisation Diagnostic for
%      Large Language Model Knowledge Graph Enrichment
%
%  Reminder from the template: no symbols, special characters, footnotes or
%  math in the title. Sub-titles are not captured in Xplore.
%-----------------------------------------------------------------------------
\title{\TODO{TITLE --- choose from the three options in the comment above}}

\author{\IEEEauthorblockN{Lynda Lagh}
\IEEEauthorblockA{\textit{TODO: department} \\
\textit{TODO: university} \\
TODO: City, Country \\
isammlyndalagh@gmail.com}
\and
\IEEEauthorblockN{TODO: Supervisor Name}
\IEEEauthorblockA{\textit{TODO: department} \\
\textit{TODO: university} \\
TODO: City, Country \\
TODO: email}
}

\maketitle

%==============================================================================
\begin{abstract}
Instruction tuning has become the standard way to adapt large language models
to knowledge graph completion, and the reported accuracies are high enough that
the task is often treated as solved. This paper asks what those accuracies
measure. We take the founding system of the subfield as an object of study
rather than as a baseline, reconstruct its pipeline stage by stage, and add the
one control it never ran: every entity name in the graph is replaced by an
opaque identifier, so that relational structure survives and surface form does
not. We report accuracy on the original and the anonymised test set as a pair,
and define the difference between them as the share of performance carried by
entity names. Three further instruments, permuted names, a seen and unseen
familiarity split, and a representation-level measure, are applied to the same
claim so that they fail in different ways. On a first dataset, ninety-one
percent of the above-chance accuracy is entity surface form. The diagnostic is
independent of any particular model and is released with the code.
\end{abstract}

\begin{IEEEkeywords}
knowledge graph completion, knowledge graph enrichment, large language models,
instruction tuning, memorisation, shortcut learning, evaluation methodology
\end{IEEEkeywords}

%==============================================================================
\section{Introduction}

Knowledge graphs are consulted as if they were complete, and the systems built
to complete them now report accuracies that appear to justify the habit.
On the WN11 benchmark, an instruction-tuned seven-billion-parameter model
reports 95.5\,\% triple classification accuracy against a chance level of 50\,\%,
rising to 96.6\,\% at thirteen billion parameters~\cite{kgllm}; comparable
systems report absolute gains of 7.2 and 7.8 mean-reciprocal-rank points over
the previous state of the art~\cite{cats,realkgc}. These numbers are produced
under an evaluation protocol in which the model is shown a triple written with
its real entity names and asked whether the triple is true.
\textbf{However, that protocol cannot distinguish a model that has learned when
a relation holds from one that has learned which names occur together}, and the
two are worth very different amounts outside the benchmark.

The distinction becomes a practical problem the moment a completion system is
used for what it is proposed for. Knowledge graph enrichment is applied where a
graph is thin: to entities that are newly added, sparsely described, or
long-tailed, and to facts that no source has yet asserted. A model that scores
a triple by recognising its entity names has, by construction, nothing to
contribute there, because the names it recognises are the names it was trained
on. Its benchmark accuracy is nonetheless high, because benchmark test entities
are drawn from the same graph as the training entities. The evaluation and the
deployment are therefore misaligned in a way that the evaluation cannot report:
the setting in which the method is measured is the one setting in which
memorisation is a winning strategy.

Existing work approaches this from two directions and stops short of it in
both. A first line adds structure to the prompt in order to improve accuracy:
latent type constraints~\cite{cats}, relation-derived soft
constraints~\cite{realkgc}, part-of-speech tags on
entities~\cite{knit}, and generated entity descriptions~\cite{ape} all raise
reported performance, and each is evaluated by comparing accuracy with the
added information against accuracy without it. Because entity names are present
in every condition of every one of these comparisons, none of them can separate
\emph{the model used the added structure} from \emph{the model used the name,
and the added tokens helped it retrieve faster}. A second line has begun to
name the underlying concern: \cite{realkgc} motivates its method by stating
that supervised fine-tuning ``inadvertently encourages models to circumvent
genuine reasoning by exploiting statistical shortcuts between graph contexts
and labels,'' and \cite{kgcf} reports that anonymising entity names in its
reasoning paths causes a significant drop, which it reads as evidence of data
leakage. The first asserts the phenomenon and measures a different one, namely
sensitivity to injected high-frequency context; the second measures the right
quantity but treats it as a robustness check on its own system rather than as a
decomposition of a reported gain. What is lost by leaving the gap open is
concrete: every accuracy figure in this subfield is currently an
unresolved mixture of two quantities with opposite implications for
deployment, and the field has no instrument that reports the ratio.

This paper supplies that instrument and applies it first to the system that
founded the subfield~\cite{kgllm}. Our contributions are:

\begin{itemize}
\item \textbf{A terminological account} that places completion inside
      enrichment rather than beside it, and that fixes the four setting
      distinctions the literature uses inconsistently: transductive versus
      inductive, closed- versus open-domain, and the closed- versus open-world
      assumption together with the measurable penalty the latter imposes
      (Section~\ref{sec:terms}).
\item \textbf{A stage-level reconstruction} of the KG-LLM pipeline as an
      eight-stage skeleton, which localises its limitation to a specific pair
      of stages rather than to the method as a whole
      (Sections~\ref{sec:skeleton}--\ref{sec:limit}).
\item \textbf{The anonymisation diagnostic}: a decomposition of above-chance
      accuracy into format compliance, surface-form memorisation, and residual
      relational knowledge, reported as a triple of numbers and never as one
      (Section~\ref{sec:method}).
\item \textbf{A seven-condition grid} that varies one factor at a time --
      names, type information, negative hardness, and negative count -- with
      the interpretation of every outcome fixed in advance
      (Section~\ref{sec:grid}).
\item \textbf{Three additional instruments} chosen because they fail in
      different ways: permuted real names, a seen and unseen familiarity split
      that leaves names intact, and a representation-level measure that,
      read alone, points the other way (Section~\ref{sec:instruments}).
\end{itemize}

We do not claim that memorisation is an error. In a transductive setting it is
a legitimate and highly effective strategy. We claim that it is not what the
literature reports itself as doing, that it is invisible under the standard
protocol, and that it does not transfer to the enrichment setting these systems
are proposed for.

%==============================================================================
\section{Terminology and Scope}
\label{sec:terms}

The vocabulary of this area is used loosely, and several of the terms that
sound like synonyms select different experiments. This section fixes each one
against a source, because the argument of the paper depends on the
distinctions.

\subsection{Knowledge graph}

Following~\cite{shi2018}, a knowledge graph is a knowledge base with at most
binary relations, written $G=(E,R,T)$ with entities $E$, relations $R$, and
triples $T=\{\langle e_i,r,e_j\rangle\}$. A \emph{surface form} is the string
that names an entity or relation; an \emph{identifier} is the symbol that
individuates it. The two are usually conflated because in every standard
benchmark they arrive together, and separating them is the mechanism of this
paper.

\subsection{Enrichment, and completion as a special case of it}

The literature treats \emph{enrichment} and \emph{completion} as near-synonyms
and, when it distinguishes them at all, does so by connotation. We adopt an
explicit containment, and the containment follows from a definition already in
print. \cite{riva2026} defines knowledge graph completion as

\begin{quote}
``the task of mining from external sources or inferring from $G$ itself new
triples $\langle e'_i, r', e'_j\rangle$ that are consistent with $G$, i.e.\ that
don't undermine its satisfiability.''
\end{quote}

Two features of this definition are unusual and both matter here. First,
\emph{mining from external sources} is inside the definition, so acquiring a
fact from outside the graph is completion and not a neighbouring activity.
Second, the success criterion is logical consistency rather than plausibility,
which places a quality condition inside the task definition itself.

Accepting that definition, however, forces a distinction it does not draw. Two
published systems both describe themselves as enriching a graph with a language
model while performing operations with entirely different risk profiles. In
\cite{sdohen}, the model adds nodes and edges: the graph gains triples, and an
incorrect output corrupts the graph. In \cite{relsem}, the model generates
concept lists attached to existing relations: no triple is added, the graph is
unchanged, and an incorrect output degrades an input string. Calling both
``enrichment'' without qualification hides the fact that only one of them can
damage $G$.

\begin{table}[htbp]
\caption{Enrichment operations. Completion is the subclass that adds triples.}
\begin{center}
\begin{tabular}{@{}p{2.05cm}p{2.5cm}p{2.6cm}@{}}
\toprule
\textbf{Operation} & \textbf{What it adds} & \textbf{Failure mode} \\
\midrule
\textbf{Completion} \newline (structural) & triples over $E$ and $R$ &
corrupts $G$: a false edge is indistinguishable from a true one \\[2pt]
\textbf{Vocabulary} \newline extension & new entities, then triples over them &
corrupts $G$, and adds spurious identifiers \\[2pt]
\textbf{Textual} \newline enrichment & descriptions of existing elements &
degrades an input string; $G$ unchanged \\[2pt]
\textbf{Schema} \newline enrichment & types, domains, ranges, inverses &
propagates: one bad constraint invalidates many triples \\[2pt]
\textbf{Confidence} \newline annotation & a score on existing triples &
mis-ranks; recoverable \\
\bottomrule
\end{tabular}
\label{tab:enrich}
\end{center}
\end{table}

We therefore use \textbf{enrichment} for any operation that adds content to
$G$, and \textbf{completion} for the subclass of enrichment operations whose
output is a triple (Table~\ref{tab:enrich}). Completion is a type of
enrichment; the reverse does not hold, since textual enrichment adds no triple.
The framing is not cosmetic. It is what makes the result of this paper a
statement about enrichment: a completion system that scores triples by
recognising entity names is precisely a system that cannot be used for the
enrichment cases -- new, thin, or long-tailed entities -- where the names carry
no learned signal.

\subsection{Task variants of completion}

Completion is evaluated through four distinct formulations, and their metrics
do not transfer between one another. A table that mixes them is misleading.

\begin{table}[htbp]
\caption{Completion task variants and their chance levels}
\begin{center}
\footnotesize
\begin{tabular}{@{}>{\raggedright\arraybackslash}p{1.7cm}%
                  >{\raggedright\arraybackslash}p{2.3cm}%
                  >{\raggedright\arraybackslash}p{1.0cm}%
                  >{\raggedright\arraybackslash}p{1.35cm}@{}}
\toprule
\textbf{Variant} & \textbf{In, out} & \textbf{Chance} &
\textbf{Used by} \\
\midrule
Triple \newline classif. & $(h,r,t)\rightarrow$ Yes/No & $0.5$ &
\cite{kgllm}, \cite{knit}, \emph{ours} \\[2pt]
Link / entity \newline prediction & $(h,r,?)\rightarrow t$ & $\approx 1/|E|$ &
\cite{cats}, \cite{realkgc}, \cite{egit} \\[2pt]
Relation \newline prediction & $(h,?,t)\rightarrow r$ & $\approx 1/|R|$ &
\cite{kgllm} \\[2pt]
Triple set \newline prediction & $G\rightarrow$ a \emph{set} & --- &
\cite{tsp} \\
\bottomrule
\end{tabular}
\label{tab:variants}
\end{center}
\end{table}

We use triple classification for the diagnostic, for one reason: its chance
level is exactly $0.5$, so the quantity we decompose -- performance above
chance -- is unambiguous. Under link prediction over tens of thousands of
candidates, chance is approximately zero and a decomposition of ``above-chance
accuracy'' has no stable denominator. Triple classification, however,
\emph{completes nothing}, which is why Section~\ref{sec:rank} converts the
classifier into a ranker and recovers link-prediction metrics without
retraining.

\subsection{Transductive and inductive settings}

A completion experiment is \textbf{transductive} when every test entity already
appears in the training graph, and \textbf{inductive} when test entities are
new while the relation set is unchanged~\cite{riva2026}. The two are frequently
used without definition. The distinction governs how much memorisation is worth:
in the transductive setting a memorised name-to-label association is directly
applicable at test time, whereas in the inductive setting it is worthless
because the test entity was never seen. The systems closest to this
work~\cite{cats,realkgc} operate inductively, which sharpens their shortcut
problem relative to ours and is precisely why our instruments are needed --
in a transductive setting, memorisation is invisible rather than penalised.

\emph{Zero-shot} is a third notion and is often conflated with inductive: it
concerns unseen \emph{supervision} for a relation, not unseen entities.
\textbf{Open-world} completion, formalised by~\cite{riva2026} after
\cite{shi2018}, expands the entity set to $E' \supset E$ while leaving $R$
unaltered, and admits triples either inferred from $T$ or extracted from an
external source; it is the inductive setting stated as a property of the task
rather than of the split.

\subsection{Closed-domain and open-domain completion}

This pair is used in two incompatible senses across the literature and, unlike
the others in this section, we could not find a source that fixes it. We
therefore state which sense we use rather than assume it.

\begin{itemize}
\item \emph{Domain coverage.} \textbf{Closed-domain} completion operates over a
      graph confined to one subject area, such as the biomedical graph
      of~\cite{sdohen}; \textbf{open-domain} completion operates over a general
      encyclopaedic graph such as Freebase or YAGO.
\item \emph{Vocabulary closure.} \textbf{Closed-domain} completion requires the
      output to be an element of the existing $E$ and $R$; \textbf{open-domain}
      completion permits new symbols to be introduced.
\end{itemize}

The second sense is the operationally consequential one, because it determines
whether a grounding step is required. A generative system that may emit any
string must map its output back to a canonical entity, and systems that omit
this step pay for it: \cite{gskgc} measures 38.9--45.3\,\% out-of-vocabulary
generation on WN18RR, and \cite{kgllm} evaluates entity prediction by substring
containment over more than one hundred thousand entities, reporting Hits@1 of
0.0949 on YAGO3-10. \textbf{In this paper, ``closed-domain'' always means
vocabulary closure}: our candidate set is the graph's own entity set, and no
new symbol is ever introduced. Our graphs are open-domain in the coverage
sense.

\subsection{Closed-world and open-world assumption}

\begin{quote}
Under the \textbf{closed-world assumption} (CWA), ``a proposition is considered
false if it is not explicitly declared true. Therefore, under the CWA, a triple
that does not exist in the KG is considered as a false triple''~\cite{tsp}.
\end{quote}

Under the \textbf{open-world assumption} (OWA), a triple absent from $G$ is
\emph{unknown} rather than false. Real graphs are incomplete, so OWA is the
honest assumption, and essentially every benchmark evaluates under CWA anyway.

The consequence is a systematic and measurable bias, which we call the
\textbf{closed-world penalty}: a correct prediction of a fact that is missing
from the graph is scored as an error. \cite{llmsim} exhibits it without naming
it, scoring its refinement two ways and finding human judgement consistently
higher than agreement with the reference graph -- 89\,\% against 82.37\,\% on
WN18RR at 10\,\% noise -- and explaining that ``the model may predict correct
triples that do not appear in the original dataset, resulting in some cases
being marked as incorrect even if they are accurate.'' \cite{ukgebn} names the
assumption; \cite{gskgc} is the first to act on it, decomposing errors into CWA
and OWA components and recovering the misjudged cases with a model jury.

The penalty bears directly on our design. Our negatives are generated by
corruption, so a triple labelled negative may be a true fact that is simply
absent from $G$. \textbf{This exposure is asymmetric across our conditions},
and the asymmetry is favourable: the anonymised arm cannot be accidentally
correct, because an opaque identifier denotes nothing outside the experiment,
whereas the real-name arm can. Any closed-world contamination therefore
inflates the real-name arm and \emph{widens} the gap we report, which means our
memorisation estimate is an upper bound and its direction of error is known.

%==============================================================================
\section{The KG-LLM Pipeline as an Eight-Stage Skeleton}
\label{sec:skeleton}

Completion was for a decade dominated by embedding methods that represent
entities and relations as vectors and score a triple by a geometric
criterion~\cite{transe}. Those methods have no surface-form component to
memorise in the sense studied here: an entity is an identifier with a learned
vector, and its name is never read. The shift to text-based completion, in
which the entity name is presented to a pretrained language model as a string,
introduced that component, and the diagnostic of this paper is meaningful only
because of it.

To locate a limitation precisely, one needs a decomposition finer than the
method as a whole. We analyse text-based completion systems as an eight-stage
skeleton, derived by reading each stage off the systems themselves rather than
by proposing an architecture, and we instantiate it here for KG-LLM~\cite{kgllm}
because that system is the object of study.

\begin{table}[htbp]
\caption{KG-LLM instantiated on the eight-stage skeleton}
\begin{center}
\begin{tabular}{@{}p{1.5cm}p{6.3cm}@{}}
\toprule
\textbf{Stage} & \textbf{KG-LLM} \\
\midrule
0 Task & triple classification, relation prediction, and entity prediction;
the first system in the area to run all three \\[2pt]
0.5 Input & text-only, static, transductive, binary, English, deterministic
with supplied negatives; entity and relation descriptions are those
of~\cite{kgbert} \\[2pt]
1 Engagement & language model as \emph{pure predictor}: no graph embedding, no
downstream scorer of any kind \\[2pt]
2 Context & $K=5$ neighbours, randomly sampled and \emph{untyped}, pipe-separated
into the prompt; $K$ is never swept \\[2pt]
2.5 Generation & \emph{skipped}: nothing is generated for the graph \\[2pt]
3 Bridge & triple linearisation into a question, plus untyped neighbourhood
verbalisation \\[2pt]
4 Reasoning & fine-tuned; LoRA~\cite{lora} with $r=8$, $\alpha=16$ on the query
and value projections; instruction tuning; no chain-of-thought, no in-context
learning \\[2pt]
5 Prediction & generative for classification and entity prediction; candidate
selection by in-prompt enumeration for relation prediction \\[2pt]
6 Resolution & \emph{nothing}: no grounding, no confidence, no calibration, no
provenance \\[2pt]
7 Evaluation & WN11, FB13, WN18RR, YAGO3-10; accuracy and Hits@1; a
$+5$\,neighbours ablation; a qualitative error table \\
\bottomrule
\end{tabular}
\label{tab:skeleton}
\end{center}
\end{table}

Two properties of Table~\ref{tab:skeleton} carry the argument of the next
section. Stage~1 is the cleanest instance of the language model as sole
predictor in the literature, which means nothing downstream can compensate for
what the model fails to learn. Stage~6 is empty, which means the system emits
no signal about its own reliability that an evaluator could inspect. Between
them, the only observable is the accuracy at Stage~7.

%==============================================================================
\section{The Limitation That Motivates This Work}
\label{sec:limit}

\subsection{Half of the founding claim was never tested}

\cite{kgllm} states that instruction tuning teaches the model ``to answer like
training triples'' \emph{and} ``to be more aware of a fact.'' The first clause
is demonstrated by the paper. The second is asserted. The headline movement --
untuned 21.1\,\% on WN11 and 9.1\,\% on FB13, rising to 89.2--96.6\,\% after
tuning -- is read by the field as fact-awareness, and \textbf{nobody has
decomposed it}.

Two features of that untuned number deserve comment, because they are the
reason the decomposition has not been done. A seven-billion-parameter model
scoring 9.1\,\% on a binary task is performing five times worse than a coin
flip, and a score far below chance is not a knowledge failure but a parsing
failure. The error table in~\cite{kgllm} confirms this directly: the untuned
failures are refusals, hedges, forum-post echoes from pretraining data, and
hallucinated personas. \emph{None of them is a wrong fact.} An eighty-point
movement measured from such a baseline is therefore not interpretable as
knowledge acquisition without a control, and no control was run.

\subsection{The standard protocol cannot see the difference}

The evaluation at Stage~7 presents each test triple with its real entity names.
A model that has learned the mapping from an entity name to a label, and a
model that has learned the conditions under which a relation holds, produce
identical outputs under this protocol whenever the test entities were seen in
training -- which, in a transductive split, is almost always. The protocol is
not weak; it is measuring a quantity that is the sum of two others and
reporting the sum.

\subsection{Negative sampling removes the pressure to do better}

KG-LLM constructs one negative per positive by uniform random corruption. A
uniformly corrupted triple is usually \emph{type-violating}, and a
type-violating negative can be rejected on the entity name alone, without
consulting the relation. Training therefore never applies pressure toward a
relational rule, because a name lookup already separates the classes. Later
systems raise the count -- six in~\cite{realkgc}, three head-corrupted and
three tail-corrupted, and twelve in~\cite{cats} -- but neither separates the
\emph{hardness} of a negative from the \emph{number} of them, so it remains
unknown which of the two is responsible for their gains.

\subsection{Statement of the gap}

Systems in this area add information to the prompt and report accuracy rising,
with entity names present in every compared condition. This design supports the
claim that the added information helps, and cannot support the claim that the
model uses it, because the two are not separated by any experiment in the
comparison. The consequence is that the reported accuracy of every such system
is an unresolved mixture of surface-form recall and relational inference, in a
proportion that has never been estimated, and the two components have opposite
value in the enrichment setting these systems are built for.

%==============================================================================
\section{Method: A Memorisation Diagnostic}
\label{sec:method}

\subsection{The decomposition}

We decompose accuracy above chance into three components:
\begin{equation}
\mathrm{acc} - c \;=\; \underbrace{\phi}_{\text{format}} \;+\;
\underbrace{\mu}_{\text{memorisation}} \;+\;
\underbrace{\kappa}_{\text{residual knowledge}}
\label{eq:decomp}
\end{equation}
with $c=0.5$ for triple classification. Each term is measured, not inferred.

\textbf{Format compliance} $\phi$ is the accuracy recovered by reading one set
of generations under a more permissive parser, i.e.\ the share of apparent
error that is an output-convention failure rather than a wrong answer. We read
each set of generations under four parsers -- the strict and the lenient rule
of~\cite{kgllm}, a constrained rule, and a generation-free comparison of the
probability assigned to the affirmative and negative tokens -- and take $\phi$
as the spread between them.

\textbf{Memorisation} $\mu$ is the \emph{gap}
\begin{equation}
\mu \;=\; \mathrm{acc}(\text{real}) - \mathrm{acc}(\text{anonymised}),
\label{eq:gap}
\end{equation}
where the anonymised arm is trained \emph{and} evaluated with every entity
surface form replaced by an opaque identifier \texttt{entity}$i$, relations
untouched:

\begin{center}
\small
\begin{tabular}{@{}ll@{}}
real: & \texttt{Is this true: dog, a domesticated} \\
      & \texttt{carnivore \_hypernym mammal?} \\[2pt]
anon: & \texttt{Is this true: entity5 \_hypernym entity0?} \\
\end{tabular}
\end{center}

Each arm is evaluated on \emph{its own} test set. Scoring an anonymised model
against real names would measure distribution shift rather than memorisation.

\textbf{Residual knowledge} $\kappa$ is the remainder: the above-chance
accuracy that survives when no name is available.

We report the three columns \emph{acc(real)}, \emph{acc(anon)}, and the gap
together, never a single accuracy. We additionally report the
\textbf{memorisation share} $\mu / (\mathrm{acc}(\text{real}) - c)$, the
fraction of above-chance performance attributable to surface form, because it
is comparable across models and datasets in a way that a raw gap is not.

\subsection{Relation to prior use of anonymisation}

Anonymisation itself is not new: \cite{kgcf} anonymises entity names in
reasoning paths and reports that ``the significant performance drop confirms''
data leakage. The difference is what the operation is used \emph{for}. There it
is a robustness ablation on the authors' own method, reported as a caution.
Here it is the measurement, applied to decompose a \emph{published} gain, with
the resulting ratio as the reported quantity. We frame the finding within
shortcut learning~\cite{geirhos2020}, the established account of decision rules
that succeed on a benchmark and fail under shifted conditions; neither
\cite{realkgc} nor \cite{cats} cites that literature despite invoking the
concept.

\subsection{The condition grid}
\label{sec:grid}

Seven training conditions vary one factor at a time
(Table~\ref{tab:grid}). Every other element of the pipeline -- data format,
prompt template, adapter rank, learning rate, token budget -- is held fixed at
the values of~\cite{kgllm}, so that any difference between rows is attributable
to the varied factor.

\begin{table}[htbp]
\caption{The condition grid. One factor changes per row.}
\begin{center}
\footnotesize
\begin{tabular}{@{}clccrl@{}}
\toprule
\textbf{ID} & \textbf{Names} & \textbf{Types} & \textbf{Neg.} &
\textbf{Inst.} & \textbf{Isolates} \\
\midrule
A & real & --- & 1 rand & 20k & baseline \\
B & anon & --- & 1 rand & 20k & entity names \\
C & anon & yes & 1 rand & 20k & types vs.\ names \\
D & anon & yes & 1 hard & 20k & neg.\ hardness \\
E & anon & yes & 6 hard & 70k & neg.\ count \\
G & real & yes & 1 rand & 20k & types w/ names \\
S & shuffled & --- & 1 rand & 20k & name binding \\
\bottomrule
\end{tabular}
\label{tab:grid}
\end{center}
\end{table}

\textbf{A} reproduces~\cite{kgllm}. \textbf{B} isolates entity names.
\textbf{C} asks whether type information can substitute for names once names
are gone; a type here is \emph{induced} from relation signatures -- the set of
relation positions an entity occupies -- rather than read from the identifier,
because the identifiers of our graphs carry no part-of-speech marker and
identifier-based extraction yields no usable types on either dataset.
Induced types are derived from structure, which anonymisation preserves, so
they are identical before and after anonymisation, which is exactly the
property condition~C requires. \textbf{D} and \textbf{E} separate the
\emph{hardness} of a negative from the \emph{count} of them, a separation that
neither \cite{realkgc} nor \cite{cats} performs. \textbf{G} is the direct
challenge to the prompt-augmentation line: if types add nothing when names are
available, the reported gains of type-aware systems are consistent with an
effect of token volume rather than of type reasoning.

A condition with twelve hard negatives, matching~\cite{cats}, was deliberately
dropped: that system is seven billion parameters and inductive, so matching its
negative count would not match its setting, and the compute is better spent
on~\textbf{G}.

\subsection{Four instruments on one claim}
\label{sec:instruments}

A single instrument invites a single objection. The objection here is specific
and fair: \emph{replacing names with opaque identifiers destroys all
information, so of course accuracy collapses}. We answer it with a measurement
rather than an argument, by applying four instruments that fail in different
ways (Table~\ref{tab:instr}).

\begin{table}[htbp]
\caption{Four instruments, four failure modes}
\begin{center}
\footnotesize
\begin{tabular}{@{}>{\raggedright\arraybackslash}p{1.85cm}%
                  >{\raggedright\arraybackslash}p{3.35cm}%
                  >{\raggedright\arraybackslash}p{1.75cm}@{}}
\toprule
\textbf{Instrument} & \textbf{What it destroys} & \textbf{Cost} \\
\midrule
Anonym.\ \newline (A to B) & all surface information &
one run \\[2pt]
Shuffled \newline names (S) & only the name-to-entity \emph{binding};
vocabulary, lengths and readability are preserved & one run \\[2pt]
Seen / unseen \newline split & \emph{nothing}; names stay intact &
free, no GPU \\[2pt]
Represent.\ \newline measure & --- (reads hidden states) & inference \\
\bottomrule
\end{tabular}
\label{tab:instr}
\end{center}
\end{table}

\textbf{Condition S} keeps every real name in the graph and permutes which
entity holds which. The permutation is a derangement, deterministic at a fixed
seed, and preserves the multiset of names exactly, so vocabulary, name lengths,
token distribution and readability are unchanged while the binding is
destroyed. If $S \approx B$, the binding was the signal and the objection is
answered by measurement. If $S \approx A$, the model never used names and
anonymisation broke something else, which would undermine the B result and must
be investigated before publication.

\textbf{The seen and unseen split} costs nothing and asks a different question.
We trained on a sample of the graph, so a substantial fraction of test entities
were never seen; splitting test accuracy by familiarity asks whether
familiarity drives accuracy while leaving every name intact. It is an
independent confirmation from the opposite direction, and the ``neither entity
seen'' bucket \emph{is} the enrichment case.

\textbf{The representation measure} is sliced mutual information between hidden
states and labels, the instrument used by the closest representation-level work
to report that fine-tuning ``primarily aligns representations rather than
injecting knowledge''~\cite{flame}. We include it because on our data it
initially points the other way, and the joint reading is itself a result
(Section~\ref{sec:disc-smi}).

\begin{figure}[htbp]
\centering
\begin{tikzpicture}[
  x=1mm, y=1mm,
  box/.style={draw, rounded corners=1.2pt, minimum width=11mm,
              minimum height=6.5mm, inner sep=0.5pt,
              font=\tiny, align=center},
  lbl/.style={font=\tiny\itshape, inner sep=1pt},
  ar/.style={-{Latex[length=1.3mm]}, thin}
]
% row 1 : the name axis
\node[box] (A) at (0,0)  {A\\real};
\node[box] (S) at (16,0) {S\\shuffled};
\node[box] (B) at (32,0) {B\\anon};
% row 2 : the constraint axis
\node[box] (C) at (32,-13) {C\\+types};
\node[box] (D) at (48,-13) {D\\+hard};
\node[box] (E) at (64,-13) {E\\+6 hard};
% the names-plus-types cell
\node[box] (G) at (0,-13) {G\\real+types};

\draw[ar] (A) -- node[above, lbl] {binding} (S);
\draw[ar] (S) -- node[above, lbl] {form}    (B);
\draw[ar] (B) -- node[right, lbl] {types}   (C);
\draw[ar] (C) -- node[above, lbl] {hard}    (D);
\draw[ar] (D) -- node[above, lbl] {count}   (E);
\draw[ar] (A) -- node[left,  lbl] {types}   (G);
\end{tikzpicture}
\caption{The grid as a ladder. Each edge removes or adds exactly one factor, so
the difference across it is attributable to that factor. The A--S--B chain
separates the name-to-entity binding from surface form in general.}
\label{fig:ladder}
\end{figure}

\subsection{Prompt ablation without retraining}

Five prompt variants are applied to the \emph{same} checkpoints: P0 bare, P1
type tags, P2 an explicit instruction, P3 both, and P4 demonstrations. P4
reproduces the mechanism of~\cite{realkgc}, whose type constraint operates by
showing other triples that use the same relation so that the model compares
head and tail types against real instances; where P1 \emph{states} the type, P4
\emph{demonstrates} it. If P4 substantially exceeds P1, demonstration beats
declaration, and the comparison is available for the price of inference.

Two cautions are recorded in advance. P2 and P3 are expected to be inert on the
tuned model, because the training loss is masked to the response and the
response is a fixed string, so instruction-following was tuned out; only the
untuned arm can react, and \emph{that contrast is itself a finding}. And
because the tuned model saw only P0, a \emph{drop} under P1--P4 may be
distribution shift, so only increases are interpretable.

\subsection{From classifier to ranker}
\label{sec:rank}

Triple classification completes nothing, and the subject of this work is
enrichment. We therefore score each candidate $t$ for a query $(h,r,?)$ by
$P(\text{Yes} \mid h,r,t)$ and sort, which yields an ordering and therefore
Hits@$K$ and mean reciprocal rank without any retraining. The principle is that
of~\cite{kgbert}, which ranks entities by a classifier score; it is not our
invention, and we use it only to make the diagnostic speak to the completion
task rather than to classification alone. This removes a
limitation that is usually stated as structural: generative decoding yields a
\emph{set} rather than an order, so mean reciprocal rank is reported as not
computable in this paradigm; scoring candidates restores it.

Ranking is 50-way and filtered, following the protocol of \cite{cats} and
\cite{realkgc}, the latter stating that it ranks each answer entity ``against
50 randomly sampled negative entities.'' Fifty-way Hits@1 is not the same
quantity as full-ranking Hits@1, and every caption in Section~\ref{sec:results}
says so.

%==============================================================================
\section{Experimental Setup}
\label{sec:setup}

\textbf{Datasets.} YAGO3-10 is primary and WN11 secondary
(Table~\ref{tab:data}). The choice is driven by the type question: YAGO3-10's
relations are strongly typed -- \texttt{wasBornIn} maps a person to a place,
\texttt{isCitizenOf} a person to a country, \texttt{playsFor} a person to a
club -- which is the structure that the type conditions require and which
WN11's lexical relations lack. On a typed graph an induced type recovers the
semantic type; on a lexical graph it does not. Running both is a decomposition
rather than a replication, since memorising on WN11 is lexical and on YAGO3-10
is world-factual.

\begin{table}[htbp]
\caption{Datasets}
\begin{center}
\footnotesize
\begin{tabular}{@{}lrrl@{}}
\toprule
 & \textbf{WN11} & \textbf{YAGO3-10} & \textbf{Note} \\
\midrule
entities & 38{,}588 & 123{,}182 & \\
relations & 11 & 37 & YAGO typed \\
test labels & shipped & generated & see below \\
typed signatures & absent & native & drives types \\
surface ambiguity & 22.7\,\% & \TODO{n} & \\
\bottomrule
\end{tabular}
\label{tab:data}
\end{center}
\end{table}

\textbf{YAGO3-10 negatives are ours, not the benchmark's.} YAGO3-10 ships no
$\pm 1$ test labels, so we generate them by corruption, filtering every
candidate against the union of the training, validation and test sets and
recording how many candidates were rejected -- that count is the closed-world
exposure of Section~\ref{sec:terms}. \textbf{Our YAGO3-10 numbers are therefore
not comparable to published YAGO3-10 results.} All comparisons are internal and
every caption states this.

\textbf{Model and adaptation.} Qwen2.5-1.5B-Instruct with
LoRA~\cite{lora} at $r=8$, $\alpha=16$, dropout 0.05, on the query and value
projections, learning rate $3\times 10^{-4}$ -- the constants of~\cite{kgllm}.
The effective batch size is 32 rather than 128, because 128 was tuned for a
corpus two orders of magnitude larger and would place warm-up at roughly a
third of training at our scale. Padding is dynamic with a 512-token cutoff.
The loss is masked to the response.

\textbf{On model scale.} A reader may reasonably suspect that a 1.5B model
memorises because it is too small to do anything else. We note that both
\cite{cats} and \cite{realkgc} report Qwen2-1.5B in their own ablations, so the
comparison is size-matched by their choice rather than by our budget. The scale
question is nonetheless live and is stated as such in
Section~\ref{sec:threats}.

\textbf{Controls.} A frozen linear probe on hidden states is included so that a
flat result under any condition is interpretable rather than ambiguous: it
distinguishes \emph{the information is absent from the representation} from
\emph{the adapter failed to extract it}. A majority-class baseline and a
degenerate-prediction check accompany every reported accuracy, because a model
that answers ``Yes'' to everything scores at the positive rate and must not be
mistaken for a working system.

\textbf{Metrics.} Accuracy, balanced accuracy, per-class precision, recall and
F1, and a confusion matrix; Hits@1/3/10 and mean reciprocal rank under the
50-way filtered protocol; expected calibration error~\cite{guo2017} and Brier
score; and a risk--coverage curve~\cite{elyaniv2010} for the abstaining
variant. The central B-versus-C contrast is run at three seeds with a paired
test and a multiple-comparison correction; a one-seed difference is not
defensible for the claim it carries.

\textbf{Reproducibility.} All conditions are declared in a single file, the
data builder asserts that any condition requesting type information actually
emits a type tag, and a validator with a non-zero exit status checks label
balance, train--test overlap, unknown identifiers, duplicates, self-loops, and
generated negatives that are in fact true. The test suite builds a
\emph{different random graph on each run} and prints its seed.

%==============================================================================
\section{Results}
\label{sec:results}

\TODO{THIS SECTION IS A SKELETON. The runs were still executing when this draft
was prepared. Conditions A and B are complete; C, D, E, G and S are pending, as
are the seen/unseen split and the representation measure. Every cell below is a
placeholder with the correct column structure. Fill from results/*.json and
delete this note.}

\subsection{Re-audit of the untuned baseline}

\begin{table}[htbp]
\caption{Untuned and tuned accuracy on \TODO{dataset}, read under four parsers.
Chance is 0.5.}
\begin{center}
\begin{tabular}{@{}lrr@{}}
\toprule
\textbf{Parser} & \textbf{Untuned} & \textbf{Tuned} \\
\midrule
strict       & \TODO{n} & \TODO{n} \\
lenient      & \TODO{n} & \TODO{n} \\
logit        & \TODO{n} & \TODO{n} \\
constrained  & \TODO{n} & \TODO{n} \\
\midrule
non-answers  & \TODO{n} & \TODO{n} \\
\bottomrule
\end{tabular}
\label{tab:parsers}
\end{center}
\end{table}

\TODO{One paragraph. State the untuned figure against KG-LLM's reported 21.1
and 9.1, and attribute the difference to the base-versus-instruction-tuned
backbone if the data support it. Give the format cost as the spread across
parsers.}

\subsection{The decomposition}

\begin{table}[htbp]
\caption{Decomposition of above-chance accuracy. Three columns, never one.}
\begin{center}
\begin{tabular}{@{}lrrrr@{}}
\toprule
\textbf{Cond.} & \textbf{acc real} & \textbf{acc anon} & \textbf{gap} &
\textbf{mem.\ share} \\
\midrule
A & \TODO{n} & \TODO{n} & \TODO{n} & \TODO{n} \\
B & \TODO{n} & \TODO{n} & \TODO{n} & \TODO{n} \\
C & \TODO{n} & \TODO{n} & \TODO{n} & \TODO{n} \\
D & \TODO{n} & \TODO{n} & \TODO{n} & \TODO{n} \\
E & \TODO{n} & \TODO{n} & \TODO{n} & \TODO{n} \\
G & \TODO{n} & \TODO{n} & \TODO{n} & \TODO{n} \\
S & \TODO{n} & \TODO{n} & \TODO{n} & \TODO{n} \\
\bottomrule
\end{tabular}
\label{tab:main}
\end{center}
\end{table}

\subsection{Corroborating instruments}

\begin{table}[htbp]
\caption{Seen and unseen familiarity split. Names are intact in every cell.}
\begin{center}
\begin{tabular}{@{}lrrr@{}}
\toprule
\textbf{Bucket} & \textbf{n} & \textbf{acc} & \textbf{mean conf.} \\
\midrule
both entities seen    & \TODO{n} & \TODO{n} & \TODO{n} \\
one seen              & \TODO{n} & \TODO{n} & \TODO{n} \\
neither seen          & \TODO{n} & \TODO{n} & \TODO{n} \\
\bottomrule
\end{tabular}
\label{tab:seenunseen}
\end{center}
\end{table}

\TODO{Report whether confidence on the ``neither seen'' bucket exceeds its
accuracy; that is overconfidence on exactly the enrichment case, and it is the
sharpest single number this table can produce.}

\subsection{Prompt ablation and recovery}

\TODO{Report as percentage of the gap recovered, not as raw accuracy. The
ceiling is the gap from Table~\ref{tab:main}. ``P4 recovers 12 percent of the
gap'' is informative; ``accuracy rose to 0.588'' is not.}

\subsection{Ranking}

\begin{table}[htbp]
\caption{Link prediction by candidate scoring. 50-way filtered, following
\cite{cats} and \cite{realkgc}; not comparable to full-ranking figures.}
\begin{center}
\begin{tabular}{@{}lrrrr@{}}
\toprule
\textbf{Cond.} & \textbf{MRR} & \textbf{H@1} & \textbf{H@3} & \textbf{H@10} \\
\midrule
A & \TODO{n} & \TODO{n} & \TODO{n} & \TODO{n} \\
B & \TODO{n} & \TODO{n} & \TODO{n} & \TODO{n} \\
\TODO{...} & & & & \\
\bottomrule
\end{tabular}
\label{tab:rank}
\end{center}
\end{table}

\subsection{Calibration and abstention}

\TODO{Expected calibration error, Brier score, and the risk--coverage curve for
each arm. State the coverage at which risk falls below the tuned model's error
rate.}

%==============================================================================
\section{Discussion}
\label{sec:disc}

\TODO{The interpretation of each outcome was fixed before the runs. Keep the
branch that occurred and delete the others.}

\subsection{What the gap means}

\TODO{If the gap is large:} a large gap means the accuracy is carried by
surface forms, and the model was made ``aware of a fact'' only in the sense of
learning an entity-name-to-label association. \TODO{If the gap is small:} the
adaptation transferred a relational signal, and the reported accuracies of this
family of systems are more robust than we anticipated -- a negative result for
our hypothesis and a positive one for the literature, which we would report as
such.

\subsection{Whether types can substitute for names}

\TODO{Branch on the C-versus-B comparison.} If $C \approx B$, type text cannot
substitute for names, and low-rank adaptation at this scale installs no usable
type rule. If $C \gg B$, the model \emph{can} use types and simply does not
when names are easier -- a more interesting result, and one that reframes the
type-augmentation line as supplying a capability the model already has rather
than teaching it one.

\subsection{Whether type information helps when names are present}

\TODO{Branch on G versus A.} If $G \approx A$, types add nothing when names are
available, which is a direct challenge to \cite{cats}, \cite{knit} and
\cite{realkgc}: their gains would then be consistent with an effect of
additional tokens rather than of type reasoning, and their module-level
ablations cannot distinguish the two because names are present in every
condition they compare. If $G > A$, types carry real signal, and the anonymised
ladder says what kind.

\subsection{Hardness against count}

\TODO{Branch on D versus C and E versus D.} If $D>C$, it is negative
\emph{hardness} that forces rule use rather than the type text, which would
imply that the gains of \cite{realkgc} and \cite{cats} owe more to their
negative sampling than to their constraint modules. If $E>D$, count matters --
but E carries 3.5 times the instances of D, so difficulty and data volume are
confounded in that row and we report instances and GPU-hours alongside it.

\subsection{The representation measure points the other way}
\label{sec:disc-smi}

\TODO{If the representation measure rises while the gap is large:} sliced
mutual information between hidden states and labels rising over the same tuning
that produces a large anonymisation gap is not a contradiction but a finding.
Representations became more label-informative, \emph{and the information they
encode is the entity name}. The measure cannot separate the two;
anonymisation can. This is the precise limitation of the closest
representation-level work~\cite{flame}, whose conclusion that fine-tuning
``primarily aligns representations rather than injecting knowledge'' rests on
that instrument alone. The result exists only because the two instruments are
reported together, which is why our evaluation emits a joint reading and flags
the case where they disagree.

\subsection{Implications}

\textbf{For evaluation.} An accuracy figure on a transductive benchmark is not
a sufficient report for a completion system. We propose that the pair
(accuracy, gap) become the minimum unit of reporting, in the same way that a
mean is not reported without a dispersion.

\textbf{For method design.} If the diagnosis holds, the productive response is
not a larger model but a training signal that makes name lookup insufficient --
which is what the hardness and count axes probe directly.

\textbf{For enrichment.} \TODO{Connect to Table~\ref{tab:enrich}: a completion
operator whose accuracy is surface form is unusable for the enrichment cases
that motivate it, because those cases are exactly the entities whose names
carry no learned signal. The ``neither seen'' bucket of
Table~\ref{tab:seenunseen} is the quantitative form of this argument.}

\textbf{On abstention.} A system that declines to answer is penalised by the
protocol of this field: \cite{kgllm}'s own error analysis shows two frontier
models declining to assert a fact about a private individual -- the
epistemically correct behaviour -- scored as wrong, because the scoring rule
reads any response containing ``not'' as a negative. Any abstaining system will
score worse on the standard benchmark than a confident guesser unless the
protocol is fixed as well, and we report risk--coverage curves for that reason.

%==============================================================================
\section{Threats to Validity}
\label{sec:threats}

\textbf{Scale.} A 1.5-billion-parameter model may memorise where a larger one
would generalise. The comparison is size-matched to the ablations
of~\cite{cats} and \cite{realkgc}, but our conclusions are stated for this
scale and a scale sweep is future work.

\textbf{Anonymisation could destroy more than names.} This is the principal
threat and the reason condition~S exists. If $S \approx A$ rather than
$S \approx B$, the B result is not safe to interpret as memorisation, and we
would report that.

\textbf{Transductive memorisation is arguably legitimate.} If the graph is
fixed, memorising is a correct strategy, and our objection would be
philosophical rather than empirical. The seen and unseen split converts the
objection into a measurement: the ``neither seen'' bucket is the case in which
memorisation cannot help, and it is present in the same test set.

\textbf{Generated negatives.} Our YAGO3-10 labels are generated, so a negative
may be an unrecorded true fact. As argued in Section~\ref{sec:terms}, this
exposure inflates the real-name arm and therefore widens the gap, making our
memorisation figure an upper bound with a known direction of error.

\textbf{Induced types are weaker than semantic types.} An induced type states
that an entity ``participates in these relations.'' We report the number of
types and their entropy rather than the count alone, because a type inventory
whose largest class covers a quarter of all entities carries less information
than its cardinality suggests.

\textbf{Prior use of anonymisation.} A survey of the corpus identified a set of
papers matching contamination-related terms, of which a subset has been
checked; all checked cases were duplicate-node leakage, privacy anonymisation,
or the inverse-relation issue rather than the measurement proposed here.
\TODO{Complete this check before asserting novelty of the instrument; state the
final counts.}

%==============================================================================
\section{Conclusion}

We asked what instruction tuning installs when a language model is adapted to
knowledge graph completion, and we answered it by adding to the founding
pipeline of the subfield the one control it never ran. Replacing entity surface
forms with opaque identifiers, while leaving relational structure intact,
splits a single accuracy figure into a memorisation term and a residual
knowledge term. On \TODO{dataset}, \TODO{n} percent of the above-chance
accuracy is entity surface form.

The instrument is independent of the system it was first applied to. It
requires one additional training run, reports a ratio rather than a score, and
can be attached to any text-based completion method without modifying it. We
also position the finding within the enrichment task it serves: completion is a
subclass of enrichment, and the enrichment cases that motivate these systems --
new, thin, and long-tailed entities -- are precisely those in which a memorised
surface form is worth nothing.

We do not claim that memorisation is a defect. We claim that it is not what
this literature reports itself as doing, that the standard protocol cannot see
it, and that reporting an accuracy without a gap is no longer sufficient.

\TODO{Future work: rationale-based supervision instead of a fixed response
string; a scale sweep; and applying the diagnostic to the inductive systems,
where the prediction is that the gap should be small and where a large one
would be a considerably stronger result.}

%==============================================================================
\section*{Acknowledgment}

\TODO{Add supervisor and institution. Put any funding in the unnumbered
footnote on the first page.}

%==============================================================================
\begin{thebibliography}{00}

\bibitem{kgllm} L. Yao, J. Peng, C. Mao, and Y. Luo, ``Exploring large language
models for knowledge graph completion,'' in \emph{Proc. IEEE Int. Conf.
Acoustics, Speech and Signal Processing (ICASSP)}, 2025. Code available at
github.com/yao8839836/kg-llm

\bibitem{kgbert} L. Yao, C. Mao, and Y. Luo, ``KG-BERT: BERT for knowledge
graph completion,'' arXiv:1909.03193, 2019.

\bibitem{transe} A. Bordes, N. Usunier, A. Garcia-Duran, J. Weston, and
O. Yakhnenko, ``Translating embeddings for modeling multi-relational data,'' in
\emph{Advances in Neural Information Processing Systems}, 2013, pp. 2787--2795.

\bibitem{lora} E. J. Hu, Y. Shen, P. Wallis, Z. Allen-Zhu, Y. Li, S. Wang,
L. Wang, and W. Chen, ``LoRA: Low-rank adaptation of large language models,''
in \emph{Int. Conf. Learning Representations (ICLR)}, 2022.

\bibitem{shi2018} B. Shi and T. Weninger, ``Open-world knowledge graph
completion,'' in \emph{Proc. AAAI Conf. Artificial Intelligence}, 2018,
pp. 1957--1964.

\bibitem{riva2026} D. Riva, A. Ferrara, and S. Montanelli, ``Open-world
knowledge graph completion with linguistic noise,'' in \emph{Proc. Int. Conf.
AI x Data and Knowledge Engineering (AIxDKE)}, 2026, pp. 106--113.

\bibitem{cats} M. Li, C. Yang, C. Xu, Z. Song, X. Jiang, J. Guo, H.-f. Leung,
and I. King, ``Context-aware inductive knowledge graph completion with latent
type constraints and subgraph reasoning,'' in \emph{Proc. AAAI Conf. Artificial
Intelligence}, 2025, p. 12106. Code available at github.com/IDEA-FinAI/CATS

\bibitem{realkgc} Y. Zhang, J. Hu, C. Xiong, Z. Qin, H. Chen, and Z. Li,
``RealKGC: Relation-constrained large language models for inductive knowledge
graph completion,'' \emph{Knowledge-Based Systems}, vol. 349, art. no. 116415,
2026. Code available at github.com/HubuKG/RealKGC

\bibitem{knit} J. Gong, X. Fang, W. Zhu, J. Peng, L. Wang, and M. Chen, ``Knit:
Toward alleviating large language model fact knowledge hallucinations in
knowledge graph completion,'' \emph{Big Data Mining and Analytics}, vol. 9,
no. 2, pp. 611--631, 2026.

\bibitem{geirhos2020} R. Geirhos, J.-H. Jacobsen, C. Michaelis, R. Zemel,
W. Brendel, M. Bethge, and F. A. Wichmann, ``Shortcut learning in deep neural
networks,'' \emph{Nature Machine Intelligence}, vol. 2, no. 11, pp. 665--673,
2020.

\bibitem{kgcf} Z. Zheng, Y. Dong, S. Wang, H. Liu, Q. Wang, and J. Li, ``KG-CF:
Knowledge graph completion with context filtering under the guidance of large
language models,'' in \emph{Proc. IEEE Int. Conf. Big Data (BigData)}, 2024,
doi: 10.1109/BigData62323.2024.10826107.

\bibitem{tsp} Y. Yuan, Y. Xu, and W. Zhang, ``Is large language model good at
triple set prediction? An empirical study,'' in \emph{Proc. IEEE Int. Conf.
Knowledge Graph (ICKG)}, 2024, pp. 470--476.

\bibitem{ukgebn} S. Yang, L. Lin, C. Liu, M. Zhang, and H. Yang, ``Prediction
of implicit relationships using uncertain knowledge graph embedding and
Bayesian networks,'' \emph{Neural Processing Letters}, vol. 57, art. 72, 2025.

\bibitem{llmsim} D. Na, N. Kertkeidkachorn, X. Liu, and K. Shirai, ``Refining
noisy knowledge graph with large language models,'' in \emph{Proc. Generative AI
and Knowledge Graph Workshop (GenAIK)}, Int. Committee on Computational
Linguistics, 2025, pp. 78--86.

\bibitem{gskgc} R. Yang, J. Zhu, J. Man, H. Liu, L. Fang, and Y. Zhou, ``GS-KGC:
A generative subgraph-based framework for knowledge graph completion with large
language models,'' arXiv:2408.10819, 2025.

\bibitem{ape} J. Zhu and P. De Meo, ``Exploring large language models for
knowledge graph completion with auto-prompting,'' in \emph{Proc. Int. Joint Conf.
Neural Networks (IJCNN)}, 2025, doi: 10.1109/IJCNN64981.2025.11228614.

\bibitem{egit} P. Zhang, X. Xu, J. Wu, X. Lu, J. Shi, X. Zhang, D. Cui,
X. Peng, S. He, P. Zong, G. Zhang, Z. Ou, M. Song, and Y. Zhu, ``Mitigating
hallucinations in knowledge graph completion via embedding-guided instruction
tuning,'' \emph{Information}, vol. 17, no. 2, art. no. 207, 2026.

\bibitem{flame} B. Xue, Y. Xu, B. Ma, Y. Song, J. Ding, L. Fu, and X. Wang,
``FLAME: Empowering frozen LLMs for knowledge graph completion,'' in \emph{Proc.
IEEE Int. Conf. Acoustics, Speech and Signal Processing (ICASSP)}, 2026,
doi: 10.1109/ICASSP55912.2026.11464010.

\bibitem{sdohen} T. Shang, S. Yang, T. Zhai, W. He, E. Mamourian, J. Zhang,
B. Hou, J. Lee, D. Duong-Tran, J. H. Moore, M. D. Ritchie, and L. Shen, ``A
novel computational analysis integrating social determinants information from
EHR and literature with Alzheimer's disease biological knowledge through large
language models and knowledge graphs,'' \emph{Innovation in Aging}, vol. 9,
no. S1, art. no. igaf102, 2025.

\bibitem{relsem} J. Si, X. Ouyang, X. Zhu, and Y. Zhang, ``A relation semantic
enhancement method for large language model based knowledge graph completion,''
in \emph{Proc. IEEE Int. Conf. Data Science in Cyberspace (DSC)}, 2024,
doi: 10.1109/DSC63484.2024.00035.

\bibitem{guo2017} C. Guo, G. Pleiss, Y. Sun, and K. Q. Weinberger, ``On
calibration of modern neural networks,'' in \emph{Proc. Int. Conf. Machine
Learning (ICML)}, 2017, pp. 1321--1330.

\bibitem{elyaniv2010} R. El-Yaniv and Y. Wiener, ``On the foundations of
noise-free selective classification,'' \emph{J. Machine Learning Research},
vol. 11, pp. 1605--1641, 2010.

\end{thebibliography}

\end{document}
